\section[Introdução]{Introdução}

O Cosmos é um site desenvolvido para fomentar uma rede colaborativa de projetos de extensão focados no acesso à justiça para populações socioeconomicamente vulneráveis. A plataforma permite que os usuários conectem diversas iniciativas, compartilhem protocolos de ação e promovam soluções conjuntas para desafios comuns. Além disso, o Cosmos facilita o acesso da comunidade a essas iniciativas, incentivando a expansão e a criação de novos projetos.

Inicialmente concebido como um aplicativo móvel, o projeto foi reformulado como um site para garantir uma execução mais eficiente e ampliar o alcance das funcionalidades essenciais.
O projeto foi idealizado por Erlane Alves dos Santos, Mestra em Ciências Criminais e Alumni PUCRS, que participou ativamente das decisões estratégicas e funcionais ao longo do desenvolvimento da plataforma.

O período do projeto ocorreu no segundo semestre de 2024 entre 5 de agosto e 4 de dezembro em conjunto do orientador Edson Ifarraguirre Moreno. 